% Werke X/1 (1917), native150–151, printed143–144; remarks by Schlesinger.
\SourceLeaf{150}\ArithmeticHead{143}{R}{EXERCITATIONES MATHEMATICAE.}
\noindent sequi: \textls[120]{sive}
\[A=B=C=D=0\]
\noindent\textls[120]{sive}
\[a'=a''=a'''=a^{\mathrm{IV}}=0.\]
\begin{center}\bfseries BEMERKUNGEN ZU DEN EXERCITATIONES MATHEMATICAE.\end{center}
\begingroup\small
\noindent\hspace*{\parindent}Zu [1.] und [2.] ist zu bemerken, daß von demselben Tage, 21. August 1796, die beiden Tagebuchaufzeichnungen
Nr. 28. und Nr. 29. herrühren; in der Nr. 28. wird auf die \textit{Exercitationes mathematicae} ausdrücklich Bezug genommen.

Die Notiz [3.] ist (nach einer schriftlichen Mitteilung von P. \textsc{Stäckel}) im Anschluß an die Beschäftigung
mit drei Abhandlungen von \textsc{Daniel Bernoulli} entstanden, die in den Novi Comment. Acad. Petrop.
\textbf{16} (1771), S. 71, \textbf{17} (1772), S. 5, \textbf{18} (1773), S. 3 enthalten sind und sich auf Reihen von der Form
\[a_1+a_2+\ldots+a_n+a_1+a_2+\ldots+a_n+a_1+a_2+\ldots\text{ in inf.}\]
\noindent beziehen, wo
\[a_1+a_2+\ldots+a_n=0\]
\noindent ist. Die ersten $n$ Glieder $a_1,a_2,\ldots,a_n$ nennt D. \textsc{Bernoulli} die \textls[120]{Periode}.\linebreak
Summiert man gliedweise und bezeichnet mit $s_\varkappa$ die Summe der $\varkappa$ ersten Reihenglieder, so ist
\[s_n=0,\ s_{n+1}=s_1,\ s_{n+2}=s_2,\ldots,s_{2n}=0,\ldots\]
\noindent und da die $n$ Summen $s_1,s_2,\ldots,s_n$ auf diese Weise gleichberechtigt sind, so schließt D. \textsc{Bernoulli}
ähnlich wie \textsc{Leibniz} bei
\[1-1+1-1+\ldots,\]
\noindent daß der wahre Wert der Summe das arithmetische Mittel aus den $n$ ersten Teilsummen ist; darin besteht
(siehe a. a. O. \textbf{18}, S. 5) das »ratiocinium metaphysicum«. —

\textsc{Gauss} setzt die Potenzreihe an
\[a_1+a_2x+a_3x^2+\ldots+a_nx^{n-1}+a_1x^n+\ldots,\]
\noindent die für $x=1$ in die betrachtete Reihe übergeht und formal aus der Entwickelung des Bruches
\[\frac{a_1+a_2x+a_3x^2+\ldots+a_nx^{n-1}}{1-x^n}\]
\noindent entsteht. Für $x=1$ ergibt sich »per regulas notas« als Grenzwert
\[-\frac{a_2+2a_3+3a_4+\ldots+(n-1)a_n}{n},\]
\noindent was wegen $s_n=0$ mit dem arithmetischen Mittel der $n$ ersten Teilsummen übereinstimmt.

Wir bemerken noch, daß D. \textsc{Bernoulli} Reihen von der hier betrachteten Art erhält, indem er in
\[\begin{gathered}
\sin x+\sin2x+\sin3x+\ldots,\\
\cos x+\cos2x+\cos3x+\ldots
\end{gathered}\]
\clearpage
\SourceLeaf{151}\ArithmeticHead{144}{L}{ANALYSIS.\quad NACHLASS.}
\noindent für $x$ einen aliquoten Teil von $\pi$ setzt. Man vergl. hierzu die Tagebuchnotiz Nr. 29. Ferner werde auf
die einschlägige neuere Literatur verwiesen\footnote{Vergl. G. \textsc{Frobenius}, \textsc{Crelles} Journal f. Mathematik \textbf{89} (1880), S. 262;
O. \textsc{Hölder}, Mathem. Annalen \textbf{20} (1882), S. 535 und zahlreiche daran anschließende Arbeiten.}.

Zu [4.] und [10.] vergl. man die Erläuterungen auf S. 18 dieses Bandes.

Die in [5.] betrachtete Umkehrung eines elliptischen Integrals kommt auch in der Notiz Nr. 32 des
\textit{Tagebuchs} vom 9. September 1796 vor und wird in der Tagebuchnotiz Nr. 33 vom 14. September verallgemeinert.
Hier sowohl wie in der Tagebuchnotiz Nr. 33 bedient sich \textsc{Gauss} einer eigentümlichen rekurrierenden
Bezeichnungsweise für die Glieder einer Reihe; es bedeutet nämlich $A$ das erste, $B$ das zweite, $C$ das dritte
Glied der betreffenden Reihe, u. s. w. Diese Bezeichnung wird von J. \textsc{Stirling}\renewcommand{\thefootnote}{**)}\footnote{J. \textsc{Stirling},
\textit{Methodus Differentialis sive Tractatus de summatione et interpolatione serierum infinitarum}, Londini 1730, S. 3.}
benutzt, der sie auf \textsc{Newton} zurückführt\renewcommand{\thefootnote}{***)}\footnote{A. a. O. S. 31.}.
Daß sich \textsc{Gauss} zu der hier in Betracht kommenden Zeit mit \textsc{Stirling} beschäftigt hat, geht auch aus anderen Zeugnissen hervor, vergl. weiter unten.

Die in [6.] mit $\Phi$ bezeichnete Funktion nannte \textsc{Gauss} später den \textls[120]{sinus lemniscaticus}; die hier
gegebene Reihenentwickelung findet sich mehrfach auch im »\textsc{Leiste}« (bei S. 20 nur bis $x^9$, bei S. 86 und
auf der Einbanddecke genau so wie hier, vergl. weiter unten) und auf dem Werke III, S. 404--406 abgedruckten
Zettel \textit{Elegantiores integralis $\displaystyle\int\frac{dx}{\surd(1-x^4)}$ proprietates}\renewcommand{\thefootnote}{$\dagger$)}\footnote{Siehe Werke III, S. 405, [4.]},
wo auch der Konvergenzbereich angegeben wird.

Die »Responsio« in [8.] läßt vermuten, daß \textsc{Gauss} durch geometrische Erwägungen veranlaßt worden
ist, die in [4.] und [10.] betrachteten zahlentheoretischen Grenzwerte zu untersuchen.

Die in [9.] betrachteten Produkte weisen auf das Studium \textsc{Euler}scher Schriften. Das erste gehört
der Theorie der elliptischen Funktionen an, vgl. den in der zweiten Abteilung dieses Bandes abgedruckten
Aufsatz des Unterzeichneten »Über \textsc{Gauss}’ Arbeiten zur Funktionentheorie« und Werke III, S. 434, 440, 464.
Das zweite
\[1-x)(1-x^2)(1-x^{2^2})(1-x^{2^3})\ldots\]
\noindent findet sich bei \textsc{Euler} in der Abhandlung \textit{De partitione numerorum} Novi Comment. Acad. Petrop. \textbf{3} (1750/51)
1753, S. 125, abgedruckt Comment. arithmeticae I, 1849, S. 73 und L. \textsc{Euleri} Opera omnia, Ser. I, vol. 2,
S. 285.

Zu [11.] vergl. man die Werke VIII, S. 30, 31 abgedruckte Notiz.
\par\noindent\hfill\textsc{Schlesinger}.\par
\ClosingRule
\clearpage
\endgroup
